\section{Smoothness of the constrained-softmax loss}
\label{sec:smoothness}

\begin{lemma}[Smoothness of $\loss(\cdot; Q)$]\label{lem:smoothness}
Under \Cref{ass:unembedding_norm}, for every $Q \in F$ the loss
$\loss(\cdot; Q): \R^d \to \R$ is $L_{\mathrm{sm}}$-smooth (i.e.\ has
$L_{\mathrm{sm}}$-Lipschitz gradient) with
\begin{align}\label{eq:smoothness_const}
   L_{\mathrm{sm}} \;\le\; \tfrac{1}{2}\, B_U^2.
\end{align}
Equivalently, $\nabla^2 \loss(x; Q) \preceq L_{\mathrm{sm}} I$
in operator norm uniformly in $x \in \R^d$.
\end{lemma}

\begin{proof}
The proof is a Hessian computation in the unembedding-image variable
$z = W_U x$, followed by the chain rule.

\noindent\textbf{Step 1: Hessian of $\loss$ in $z$-coordinates.}
Define $\ell: \R^{|\mathcal V|} \to \R$ by
$\ell(z) \coloneqq -\log\!\bigl(\1_{\Cset(Q)}^\top \softmax(z)\bigr)$, so
that $\loss(x; Q) = \ell(W_U x)$. By the same softmax-gradient calculation
of the proof of \Cref{lem:gradient_form} applied at the
$|\mathcal V|$-dimensional level (replacing $W_U$ by the identity),
\begin{align*}
   \nabla \ell(z) \;=\; p(z) \;-\; \qcond(z),
\end{align*}
where $p(z) \coloneqq \softmax(z)$ and
$\qcond(z)_v \coloneqq p(z)_v \1[v \in \Cset(Q)] / (\1_{\Cset(Q)}^\top p(z))$
are the $\R^{|\mathcal V|}$-versions of \Cref{eq:z_p_def,eq:qcond_def}.
Differentiating once more in $z$, the Jacobians of the two softmax-type
maps are
\begin{align*}
   \nabla p(z) \;=\; \mathrm{diag}(p) \;-\; p p^\top,
   \qquad
   \nabla \qcond(z) \;=\; \mathrm{diag}(\qcond) \;-\; \qcond\, \qcond^\top,
\end{align*}
(both standard softmax-Jacobian identities: the second follows from the
first applied to the renormalised distribution $\qcond$, which is itself
the softmax of $z + \log(\1_{\Cset(Q)}) = z$ restricted to the support
$\Cset(Q)$, and treating $\qcond$ as a distribution on all of $\mathcal V$
with zeros off $\Cset(Q)$). Therefore
\begin{align}\label{eq:Hess_ell}
   \nabla^2 \ell(z)
   \;=\; \bigl(\mathrm{diag}(p) - p p^\top\bigr)
       \;-\; \bigl(\mathrm{diag}(\qcond) - \qcond\, \qcond^\top\bigr).
\end{align}

\noindent\textbf{Step 2: operator-norm bound on $\nabla^2 \ell(z)$.}
Each summand in \Cref{eq:Hess_ell} is the covariance matrix of a
$\mathcal V$-valued random variable (with distribution $p$ and
$\qcond$ respectively), and therefore positive semi-definite. The
operator norm of a covariance matrix
$\mathrm{Cov}(X) = \mathrm{diag}(p) - p p^\top$ for $X \in \mathcal V$
with distribution $p$ is at most $\max_v p_v (1 - p_v) \le 1/4$
(since $t(1-t) \le 1/4$ for $t \in [0, 1]$). Hence
$\norm{\mathrm{diag}(p) - p p^\top}_{\mathrm{op}} \le 1/4$
and likewise
$\norm{\mathrm{diag}(\qcond) - \qcond\, \qcond^\top}_{\mathrm{op}} \le 1/4$.
The triangle inequality gives
\begin{align}\label{eq:Hess_ell_op_bound}
   \norm{\nabla^2 \ell(z)}_{\mathrm{op}}
   \;\le\; \tfrac{1}{4} \;+\; \tfrac{1}{4}
   \;=\; \tfrac{1}{2}.
\end{align}

\noindent\textbf{Step 3: chain rule back to $\loss$.}
Since $\loss(x; Q) = \ell(W_U x)$, the chain rule for second derivatives
gives $\nabla^2 \loss(x; Q) = W_U^\top \nabla^2 \ell(W_U x)\, W_U$, and
therefore
\begin{align*}
   \norm{\nabla^2 \loss(x; Q)}_{\mathrm{op}}
   \;\le\; \norm{W_U}_{\mathrm{op}}^2 \cdot \norm{\nabla^2 \ell(W_U x)}_{\mathrm{op}}
   \;\le\; B_U^2 \cdot \tfrac{1}{2}
   \;=\; \tfrac{1}{2}\, B_U^2,
\end{align*}
where the second inequality uses \Cref{ass:unembedding_norm} and
\Cref{eq:Hess_ell_op_bound}. This bound holds uniformly in $x$ and
$Q$, which is the assertion.
\end{proof}

\begin{remark}\label{rem:smoothness_non_convex}
\Cref{lem:smoothness} gives smoothness of $\loss$ but does \emph{not}
imply convexity. The Hessian decomposition \Cref{eq:Hess_ell} subtracts
the $\qcond$-covariance (a PSD matrix) from the $p$-covariance (also
PSD), and the difference is in general \emph{not} PSD: when the
$\qcond$-covariance dominates, $\nabla^2 \ell(z) \prec 0$ along some
directions, so $\loss$ is non-convex. This is intrinsic to the
constrained-softmax loss: maximising a sum of softmax components is a
non-convex selection problem~\cite{bottou2018optimization}. The
non-convexity has two consequences for our analysis: (i) we use the
\emph{descent-lemma} route through smoothness alone, not a
convexity-based potential argument, and (ii) the headline rate of
\Cref{thm:main_convergence_biased_sgd} is the standard non-convex SGD
rate $O(1/\sqrt T)$ to a stationary point with bias floor, not the faster
$O(1/T)$ that would follow from convexity. We discuss the
implications for empirical practice in \Cref{sec:discussion}.
\end{remark}

\begin{remark}\label{rem:smoothness_constant_tightness}
The constant $1/2$ in $L_{\mathrm{sm}} \le \tfrac{1}{2} B_U^2$ comes from
the operator-norm bound $\norm{\nabla^2 \ell(z)}_{\mathrm{op}} \le 1/2$
via the triangle inequality. A tighter analysis, using that the two
covariance matrices are not independent (the $\qcond$ is a renormalised
restriction of $p$), can in principle improve the constant, but
$L_{\mathrm{sm}} \le \tfrac{1}{2} B_U^2$ is sufficient for the
qualitative rate analysis below and pins no constants beyond
$B_U$ from \Cref{ass:unembedding_norm}.
\end{remark}
